\documentclass[11pt,a4paper]{article}

\usepackage{amsmath,amssymb,amsthm}
\usepackage{mathtools}
\usepackage{hyperref}
\usepackage{geometry}
\usepackage{enumitem}
\usepackage{tikz}
\usepackage{tikz-cd}
\usetikzlibrary{arrows.meta,calc,positioning}

\geometry{margin=1in}

\theoremstyle{definition}
\newtheorem{definition}{Definition}[section]
\newtheorem{example}[definition]{Example}
\newtheorem{remark}[definition]{Remark}
\newtheorem{construction}[definition]{Construction}
\newtheorem{notation}[definition]{Notation}
\newtheorem{question}[definition]{Question}

\theoremstyle{plain}
\newtheorem{theorem}[definition]{Theorem}
\newtheorem{lemma}[definition]{Lemma}
\newtheorem{proposition}[definition]{Proposition}
\newtheorem{corollary}[definition]{Corollary}

\newcommand{\R}{\mathbb{R}}
\newcommand{\Z}{\mathbb{Z}}
\newcommand{\F}{\mathbb{F}}
\newcommand{\cF}{\mathcal{F}}
\newcommand{\cE}{\mathcal{E}}
\newcommand{\cB}{\mathcal{B}}
\newcommand{\Prob}{\mathcal{P}}

\title{Embedding the MPC Ambient Diagram into Structured Categories}
\author{Research Notes}
\date{January 4, 2026}

\begin{document}

\maketitle

\begin{abstract}
This note treats the standard ``MPC ambient diagram'' as a \emph{presented diagram shape} (a small category generated by named nodes/arrows together with commutativity relations) and studies its \emph{interpretations} as functors into other categories.
The point is to separate (i) the \emph{syntax} of the arrow-translation template from (ii) the choice of \emph{semantic category} in which the template is installed.
We outline a theorem-style framework for interpreting the same diagram in cartesian categories, in stochastic/Markov-kernel settings, in quantum-channel categories, and in geometric categories such as $\mathbf{Top}$ and $\mathbf{Man}$.
Each target section gives an explicit node-by-node and arrow-by-arrow reading of the objects (execution, inputs, shares/transcripts, views, outputs) and morphisms (projections, protocol maps, view restrictions, reconstruction) in the ambient diagram.
\end{abstract}

\tableofcontents

\section{Introduction}

Multi-party computation (MPC) enables parties to compute a function of their private inputs while controlling what information is revealed through an adversary's \emph{view}.
This note focuses on the \emph{diagrammatic} skeleton used to describe a protocol:
one identifies an ``execution object'' $E$ (inputs plus randomness or other internal degrees of freedom) and records a fixed family of named arrows out of $E$ (share/transcript maps, output maps, and adversarial view maps).

The goal here is not to commit to any single base category (nor to any ``finite/discrete'' baseline), but rather to treat the ambient diagram as a \emph{pure categorical pattern} that can be installed into different semantic categories.
This provides a controlled way to research ``MPC-like'' mathematical structures: we can ask which categorical assumptions are required to even state the diagram, and how the meaning of each arrow changes across settings (deterministic, stochastic, quantum, topological, smooth).

\subsection{Fundamental MPC View}

We begin with the correctness viewpoint of MPC, without introducing any adversary.
Fix a category $\mathcal{C}$ in which we will interpret the diagram.
At minimum, $\mathcal{C}$ must have the products appearing in Figure~\ref{fig:protocol-correctness} so that we can speak of objects like $\prod_i X_i$ and $\prod_i S_i$ and the usual coordinate projections.

The space of executions is an object $E$ equipped with a map $\pi_X:E\to \prod_i X_i$ that forgets internal randomness/state.
The intended functionality is a morphism $f:\prod_i X_i\to Y$; hence the ``ideal'' output map from executions is the composite $f\circ \pi_X:E\to Y$.
A concrete protocol supplies another morphism $\pi_S:E\to \prod_i S_i$ (full transcript/share vector) together with a reconstruction morphism $\rho:\prod_i S_i\to Y$.
Protocol correctness is exactly the statement that these two parallel arrows $E\rightrightarrows Y$ coincide, i.e.\ the square in Figure~\ref{fig:protocol-correctness} commutes:
\[
\rho\circ \pi_S \;=\; f\circ \pi_X.
\]
Equivalently, $\rho$ exhibits $f\circ \pi_X$ as a factorization through $\pi_S$ (no uniqueness is claimed).

\begin{notation}
We fix the following \emph{operational names} for nodes and arrows in the template.
In concrete algebraic protocols one often chooses $X_i\subseteq \F_q^{d_i}$ and $Y=\F_q^k$, but in this note $X_i,Y,\Omega_i,\dots$ are treated abstractly as objects of~$\mathcal{C}$.
\begin{itemize}
    \item $X_i$: the input object of party $i$, for $i \in \{1, \ldots, n\}$.
    \item $\Omega_i$: an ``internal randomness/state'' object for party $i$.
    \item $\Omega := \prod_{i=1}^n \Omega_i$: the total internal randomness/state object.
    \item $E$: the execution object. In many semantics one models $E\cong (\prod_i X_i)\times \Omega$, but the template only uses that $E$ comes with named observation maps (e.g.\ $\pi_X,\pi_S$).
    \item $\pi_X: E \to \prod_{i=1}^n X_i$: the map forgetting internal randomness/state.
    \item $Y$: the output object.
    \item $f: \prod_{i=1}^n X_i \to Y$: the functionality.
\end{itemize}
\end{notation}

\begin{definition}[Protocol Map]
\label{def:protocol-map-embeddings-note}
An $n$-party protocol for computing $f$ in a base category $\mathcal{C}$ consists of:
\begin{enumerate}
    \item An object $S_i$ for each party $i$ (``share/transcript object'').
    \item A protocol map $\pi_S: E \to \prod_{i=1}^n S_i$ assigning shares/transcripts to parties.
    \item A reconstruction map $\rho: \prod_{i=1}^n S_i \to Y$ such that the diagram in Figure~\ref{fig:protocol-correctness} commutes, i.e.\ $\rho \circ \pi_S = f \circ \pi_X$.
\end{enumerate}
\end{definition}

\begin{figure}[ht]
\centering
\begin{tikzcd}
E \arrow[r, "\pi_S"] \arrow[d, "\pi_X"'] & \prod_{i=1}^n S_i \arrow[d, "\rho"] \\
\prod_{i=1}^n X_i \arrow[r, "f"'] & Y
\end{tikzcd}
\caption{Correctness as a commuting square in a base category $\mathcal{C}$ with the relevant products: the two composites $E \xrightarrow{\pi_S} \prod_i S_i \xrightarrow{\rho} Y$ and $E \xrightarrow{\pi_X} \prod_i X_i \xrightarrow{f} Y$ coincide.}
\label{fig:protocol-correctness}
\end{figure}

\subsection{Correctness via auxiliary elimination: three categorical mechanisms}
\label{subsec:combine-cancel}
Ignoring adversarial views, the correctness square $\rho\circ \pi_S=f\circ \pi_X$ can be read as: a protocol map $\pi_S$ \emph{combines} inputs with auxiliary resources (often denoted $\omega_i$), and reconstruction $\rho$ \emph{cancels} (or quotients out) the auxiliary contribution so that the net effect agrees with the ideal functionality.
This subsection records three common categorical patterns that make this \emph{auxiliary-elimination} intuition mathematically explicit.
\emph{Scope:} these are correctness-side templates only; privacy/leakage requires reintroducing adversarial view arrows in the ambient diagram.

\subsubsection{(A) Group-object actions and torsor masking (classical additive-style)}
\label{subsubsec:combine-cancel-group}
Assume $\mathcal{C}$ has finite products and admits (internal) group objects.
Let $G$ be a group object.
To keep types explicit, we assume each share/transcript object $S_i$ carries a right $G$-action and we have a structure map
\[
u_i:\; X_i \longrightarrow S_i,
\]
so that the payload $x_i\in X_i$ is first placed inside the share space and then masked by the $G$-action.
Concretely, fix action morphisms
\[
\beta_i:\; S_i\times G \longrightarrow S_i,
\qquad
(s_i,g)\longmapsto s_i\cdot g.
\]
Let $\omega_i$ be modeled as an element of $G$ (or, more generally, as a point of a $G$-torsor).
Then a basic masking morphism is
\[
\mathrm{mask}_i:\; X_i\times G \longrightarrow S_i,
\qquad
\mathrm{mask}_i(x_i,\omega_i):=u_i(x_i)\cdot \omega_i
\quad\text{(in additive notation: }u_i(x_i)+\omega_i\text{)}.
\]
At the $n$-party level, with $X:=\prod_i X_i$ and $\Omega:=G^n$, one can take $E=X\times \Omega$ and define
\[
\pi_S:=\prod_{i=1}^n \mathrm{mask}_i:\; X\times G^n \longrightarrow \prod_i S_i.
\]
``Cancellation'' is then expressed by requiring $\rho$ to be invariant along the $G$-orbits generated by the masks (or to use algebraic relations among the $\omega_i$ so that the total mask contribution is neutral).
In the simplest additive secret-sharing pattern, one takes $G$ abelian and restricts the auxiliary object to masks whose total product is the identity.
Categorically, rather than taking $\Omega=G^n$, define $\Omega\hookrightarrow G^n$ as the fiber over the identity of the multiplication morphism
\[
\mu:\; G^n \longrightarrow G,
\qquad
(g_1,\dots,g_n)\longmapsto g_1\cdots g_n,
\]
equivalently as the pullback/equalizer of $\mu$ against the unit $e:1\to G$.
In element notation, $(\omega_1,\dots,\omega_n)\in\Omega$ satisfies
\[
\omega_1\cdot \omega_2\cdots \omega_n = e,
\]
so that reconstruction can cancel the total mask contribution.
One categorical way to express cancellation is: $\rho$ factors through the coinvariants/coequalizer of the induced $G$-action on $\prod_i S_i$ (hence is constant on masking orbits).

\subsubsection{(B) Monoidal cancellation via duals (evaluation/coevaluation)}
\label{subsubsec:combine-cancel-duals}
Assume $\mathcal{C}$ is symmetric monoidal $(\mathcal{C},\otimes,I)$ and \emph{rigid} (every object has a left/right dual).
``Auxiliary resources'' can be modeled by introducing, for each party $i$, a dual pair $(W_i,W_i^\vee)$ with coevaluation and evaluation maps
\[
\mathrm{coev}_i:\; I \longrightarrow W_i\otimes W_i^\vee,
\qquad
\mathrm{ev}_i:\; W_i^\vee\otimes W_i \longrightarrow I,
\]
satisfying the zig--zag (yanking) identities
\[
(\mathrm{id}_{W_i}\otimes \mathrm{ev}_i)\circ (\mathrm{coev}_i\otimes \mathrm{id}_{W_i})=\mathrm{id}_{W_i},
\qquad
(\mathrm{ev}_i\otimes \mathrm{id}_{W_i^\vee})\circ (\mathrm{id}_{W_i^\vee}\otimes \mathrm{coev}_i)=\mathrm{id}_{W_i^\vee}.
\]
An abstract way to formalize ``combine'' is to let the encoding insert a dual pair (or correlate workspaces) and then couple inputs to one side of the pair.
For example, an encoding can have the schematic form
\[
X_i \cong X_i\otimes I
\xrightarrow{\ \mathrm{id}\otimes \mathrm{coev}_i\ }
X_i\otimes W_i\otimes W_i^\vee
\xrightarrow{\ \mathrm{enc}_i\ }
S_i,
\]
where $\mathrm{enc}_i$ is the protocol-specific interaction between the input register and an auxiliary dual pair.
``Cancellation'' is provided when reconstruction composes with $\mathrm{ev}_i$ in a way that matches the inserted $\mathrm{coev}_i$:
the point is not that $\mathrm{ev}_i$ cancels arbitrary ancillas, but that $\mathrm{coev}_i$ followed by $\mathrm{ev}_i$ cancels by the yanking identities when wired compatibly.
This pattern is the categorical backbone of many quantum/process-theoretic ``introduce ancilla then discard'' constructions: the cancellation is not a group inverse but the evaluation/cap-cup cancellation arising from duality.

\subsubsection{(C) Cancellation as quotienting: coequalizers and invariants}
\label{subsubsec:combine-cancel-quotient}
Assume $\mathcal{C}$ has the relevant colimits (in particular coequalizers).
Let $\pi_S:E\to S$ be the protocol map (here $S:=\prod_i S_i$).
To say ``$\omega$ does not matter to the decoded output'' can be expressed by identifying an equivalence relation on $S$ generated by varying $\omega$ while keeping the underlying input fixed.
Abstractly, suppose there is a ``masking'' endomorphism (or family of endomorphisms) $m:S\rightrightarrows S$ induced by changing auxiliary degrees of freedom, and form its coequalizer
\[
S \xrightarrow{q} S/{\sim}
\]
so that $q\circ m_1=q\circ m_2$.
Then a cancellation requirement is exactly the existence of a factorization
\[
\rho:\; S \longrightarrow Y
\qquad\text{with}\qquad
\rho=\overline{\rho}\circ q
\]
for some $\overline{\rho}:S/{\sim}\to Y$.
Equivalently: $\rho$ is constant on the identified ``masking'' congruence classes.
There is also a clean formulation directly on the execution side.
Let $X:=\prod_i X_i$ and suppose $\pi_X:E\to X$ presents the ``forget auxiliary'' quotient in the sense that $\pi_X$ is (or is treated as) a coequalizer of its kernel pair:
\[
E\times_X E \rightrightarrows E \xrightarrow{\ \pi_X\ } X.
\]
Then any morphism $h:E\to Y$ that equalizes the two projections $E\times_X E\rightrightarrows E$ factors uniquely through $\pi_X$:
there exists a unique $\bar h:X\to Y$ with $h=\bar h\circ \pi_X$.
Applying this with $h:=\rho\circ \pi_S$ yields $\rho\circ\pi_S=\bar h\circ \pi_X$; correctness is the additional identification $\bar h=f$.

\section{Ambient Diagram}
\label{sec:ambient-diagram-embeddings-note}

Correctness alone concerns commuting paths from the execution object to outputs (Figure~\ref{fig:protocol-correctness}).
For security analysis and for cross-protocol comparison, we extend the basic correctness square by adjoining explicit \emph{view/observation} morphisms out of the same execution object~$E$.
The resulting ambient diagram (Figure~\ref{fig:master-diagram}) is the arrow-translation template: protocols are compared by instantiating the same named nodes/arrows and then interpreting what each arrow means in the chosen semantic category.

\begin{figure}[ht]
\centering
\begin{tikzcd}[row sep={3.0em,between origins}, column sep=huge]
& E
  \arrow[r, "\pi_S"]
  \arrow[d, "\pi_X"']
  \arrow[dl, "\pi_{X_H}"'{pos=0.45}]
  \arrow[rr, bend left=18, "\pi_{\mathrm{full}}"{pos=0.55}]
  \arrow[rrr, bend left=22, "\pi_A"{pos=0.55}]
  \arrow[dr, "\pi_{\mathrm{out}}"{pos=0.35, bend right=20}]
& \prod_{i=1}^n S_i \arrow[d, "\rho"]
& B_{\mathrm{full}}
  \arrow[l, "\mathrm{proj}_S"']
  \arrow[r, "\mathrm{proj}_A"]
  \arrow[dl, "\rho\,\circ\,\mathrm{proj}_S"{pos=0.6}]
& B_A \arrow[l, "\iota_A"', hook, bend left=35]
\\
X_H
& \prod_{i=1}^n X_i \arrow[l, "\mathrm{proj}_{X_H}"] \arrow[r, "f"']
& Y
&
&
\end{tikzcd}
\caption{Ambient diagram (template). The top-left square is the correctness square extended by a maximal transcript/view object $B_{\mathrm{full}}$ and an adversarial view object $B_A$, together with post-processing arrows $\mathrm{proj}_S$ and $\mathrm{proj}_A$.}
\label{fig:master-diagram}
\end{figure}

\subsection{How to read the ambient diagram}
\label{subsec:ambient-diagram-legend-embeddings-note}

\textbf{Adversarial conventions.}
\begin{itemize}
    \item $A \subseteq \{1, \ldots, n\}$: the set of corrupted parties, and $H = \{1, \ldots, n\} \setminus A$ the set of honest parties.
    \item $X_H := \prod_{i \in H} X_i$ and $X_A := \prod_{i \in A} X_i$, with $\mathrm{proj}_{X_H}: \prod_{i=1}^n X_i \to X_H$ the projection forgetting adversarial inputs.
\end{itemize}

\textbf{View objects.}
\begin{itemize}
    \item $B_{\mathrm{full}}$: a ``maximal transcript'' object used as a common ambient for comparing smaller views by post-processing.
    \item $B_A$: an adversarial view object associated to corruption set~$A$.
\end{itemize}

\textbf{View maps (observation morphisms).}
\begin{itemize}
    \item $\pi_{\mathrm{full}}: E \to B_{\mathrm{full}}$: a maximal transcript/observation map.
    \item $\pi_A:E\to B_A$: an adversarial view map obtained from $\pi_{\mathrm{full}}$ by post-processing.
\end{itemize}

\textbf{Post-processing maps.}
\begin{itemize}
    \item $\mathrm{proj}_S: B_{\mathrm{full}} \to \prod_i S_i$: projection/post-processing to shares.
    \item $\mathrm{proj}_A: B_{\mathrm{full}} \to B_A$: post-processing to the adversary's view.
    \item $\iota_A: B_A \hookrightarrow B_{\mathrm{full}}$: a padding inclusion used for bookkeeping comparisons inside a common ambient object (typically non-operational; often non-canonical unless additional basepoint/section data are chosen).
\end{itemize}

\textbf{Commutativity relations (as read off from Figure~\ref{fig:master-diagram}).}
\begin{enumerate}
    \item $\rho \circ \pi_S = f \circ \pi_X$ (protocol correctness).
    \item $\pi_{\mathrm{out}} = f \circ \pi_X = \rho \circ \pi_S$ (output as either ideal functionality or reconstruction).
    \item $\pi_A = \mathrm{proj}_A \circ \pi_{\mathrm{full}}$ (adversarial view as a post-processing of the full transcript).
\end{enumerate}

\paragraph{Information lens (factorization preorder induced by post-processing).}
Even without committing to a specific semantic category, there is a canonical \emph{preorder} on observations out of $E$ induced by post-processing:
given two arrows $\pi_1:E\to B_1$ and $\pi_2:E\to B_2$ in $\mathcal{C}$, we say that $\pi_2$ is \emph{at least as informative as} $\pi_1$ if there exists a morphism $q:B_2\to B_1$ with $\pi_1=q\circ \pi_2$.
In the ambient diagram, $\pi_{\mathrm{full}}$ is at least as informative as $\pi_A$ because $\pi_A=\mathrm{proj}_A\circ \pi_{\mathrm{full}}$.

\section{A General Framework: Presented Diagrams and Functorial Interpretations}
\label{sec:general-framework}

This section formalizes the sense in which the ambient diagram is a purely syntactic object that can be installed into different semantic categories.

\subsection{The presented ambient-diagram category}
\label{subsec:presented-category}
Let $\mathsf{AD}$ denote the small category \emph{presented} by:
\begin{itemize}
    \item objects: the node symbols appearing in Figure~\ref{fig:master-diagram} (e.g.\ $E$, $\prod_i X_i$, $\prod_i S_i$, $B_{\mathrm{full}}$, $B_A$, $Y$, $X_H$);
    \item generating morphisms: the named arrows (e.g.\ $\pi_X,\pi_S,\rho,f,\pi_{\mathrm{full}},\pi_A,\mathrm{proj}_A,\mathrm{proj}_S,\mathrm{proj}_{X_H},\pi_{X_H},\pi_{\mathrm{out}},\iota_A$);
    \item relations: the commutativity constraints listed in Section~\ref{subsec:ambient-diagram-legend-embeddings-note}.
\end{itemize}
Equivalently, $\mathsf{AD}$ is the free category on the ambient-diagram directed graph modulo the stated relations.

\subsection{Interpretations as functors}
\begin{definition}[Interpretation / installation of the ambient diagram]
\label{def:interpretation}
Let $\mathcal{C}$ be a target category.
An \textbf{interpretation} (or \textbf{installation}) of the ambient diagram in $\mathcal{C}$ is a functor
\[
F:\mathsf{AD}\longrightarrow \mathcal{C}.
\]
It assigns objects to nodes and morphisms to arrows in a way that preserves composition and enforces the defining relations.
\end{definition}

\begin{remark}[Structural requirements live in the target category]
\label{rem:structure-in-target}
The ambient diagram itself does not determine what ``product,'' ``projection,'' ``(co)fibration,'' or ``information'' should mean.
Those notions only become available after selecting a target category $\mathcal{C}$ and (when relevant) extra structure on $\mathcal{C}$ (e.g.\ finite products, a symmetric monoidal structure, a Markov structure, or a Quillen model structure).
\end{remark}

\subsection{Arrow-by-arrow schema (the abstract checklist)}
\label{subsec:abstract-checklist}
For later sections we record the abstract node/arrow list that must be interpreted.\footnote{In practice one often omits the padding inclusion $\iota_A$ unless a canonical padding choice exists in the target semantics.}
\begin{description}
    \item[\textbf{Nodes.}]
    $E$, $\prod_i X_i$, $X_H$, $\prod_i S_i$, $Y$, $B_{\mathrm{full}}$, $B_A$.
    \item[\textbf{Core arrows (correctness square).}]
    $\pi_X:E\to \prod_i X_i$, $\pi_S:E\to \prod_i S_i$, $\rho:\prod_i S_i\to Y$, $f:\prod_i X_i\to Y$.
    \item[\textbf{Adversary/view arrows.}]
    $\pi_{\mathrm{full}}:E\to B_{\mathrm{full}}$, $\mathrm{proj}_S:B_{\mathrm{full}}\to \prod_i S_i$, $\mathrm{proj}_A:B_{\mathrm{full}}\to B_A$, $\pi_A:E\to B_A$.
    \item[\textbf{Input split arrow.}]
    $\mathrm{proj}_{X_H}:\prod_i X_i\to X_H$ and $\pi_{X_H}:E\to X_H$.
    \item[\textbf{Output arrow.}]
    $\pi_{\mathrm{out}}:E\to Y$.
    \item[\textbf{Bookkeeping arrow (optional).}]
    $\iota_A:B_A\hookrightarrow B_{\mathrm{full}}$.
\end{description}

\section{Embedding into Markov-Kernel Categories}
\label{sec:markov}

This section interprets the ambient diagram in a stochastic setting where morphisms represent Markov kernels (conditional distributions) rather than deterministic functions.

\subsection{Target category and basic axioms}
We fix a category $\mathcal{C}_{\mathrm{stoch}}$ of measurable spaces and Markov kernels (often called $\mathbf{Stoch}$).
Objects are measurable spaces; morphisms $X\leadsto Y$ are Markov kernels $K(\,\cdot\,\mid x)$ from $X$ to probability measures on $Y$.
Deterministic measurable maps $g:X\to Y$ embed as deterministic kernels $x\mapsto \delta_{g(x)}$.
Composition is Chapman--Kolmogorov (integration against kernels).

\begin{remark}[What changes vs.\ deterministic semantics]
In a kernel category, the protocol map $\pi_S$ and view map $\pi_A$ are naturally modeled as \emph{kernels} rather than functions.
The commutativity constraint ``$\rho\circ\pi_S=f\circ \pi_X$'' must therefore be interpreted either as equality of kernels (strong, pointwise) or as equality after pushing forward a chosen prior (weaker, ``almost sure'').
\end{remark}

\subsection{Node-by-node interpretation (faithful to Figure~\ref{fig:master-diagram})}
\begin{description}
    \item[\textbf{$X_i$.}] Measurable input space of party $i$.
    \item[\textbf{$\prod_i X_i$.}] Product measurable space of all inputs (cartesian product with the product $\sigma$-algebra).
    \item[\textbf{$A,H,X_H$.}] Fix $A\subseteq\{1,\dots,n\}$ and $H=\{1,\dots,n\}\setminus A$; then $X_H=\prod_{i\in H}X_i$ as a measurable space.
    \item[\textbf{$\Omega$.}] Measurable ``internal randomness/state'' space (if modeled explicitly).
    \item[\textbf{$E$.}] Execution measurable space. A canonical choice is $E=(\prod_i X_i)\times \Omega$; alternatively one can take $E=\prod_i X_i$ and absorb the randomness into kernels $\pi_S,\pi_{\mathrm{full}}$.
    \item[\textbf{$S_i$.}] Measurable share/transcript space for party $i$.
    \item[\textbf{$\prod_i S_i$.}] Product measurable space of the full share vector / full transcript bundle.
    \item[\textbf{$Y$.}] Measurable output space.
    \item[\textbf{$B_{\mathrm{full}}$.}] Measurable ``maximal transcript'' space (a common ambient observation alphabet).
    \item[\textbf{$B_A$.}] Measurable adversarial-view alphabet for corruption set $A$.
\end{description}

\subsection{Arrow-by-arrow interpretation (every arrow in Figure~\ref{fig:master-diagram})}
\begin{description}
    \item[\textbf{$\pi_X:E\to \prod_i X_i$.}]
    If $E=(\prod_i X_i)\times \Omega$, take $\pi_X$ to be the deterministic projection $(\chi,\omega)\mapsto \chi$ (embedded as a deterministic kernel).
    If $E=\prod_i X_i$, take $\pi_X=\mathrm{id}$.

    \item[\textbf{$f:\prod_i X_i\to Y$.}]
    Deterministic measurable functionality map (again a deterministic kernel).

    \item[\textbf{$\pi_S:E\leadsto \prod_i S_i$.}]
    The protocol channel (kernel) producing the full share/transcript vector from the execution state.
    In the explicit-randomness model $E=(\prod_i X_i)\times\Omega$, this may be deterministic as a function of $(\chi,\omega)$; in the kernel-only model it is genuinely stochastic.

    \item[\textbf{$\rho:\prod_i S_i\to Y$.}]
    Deterministic measurable reconstruction/post-processing map.

    \item[\textbf{$\pi_{\mathrm{out}}:E\leadsto Y$.}]
    The \emph{output observation kernel}. In a deterministic correctness model one sets
    \[
    \pi_{\mathrm{out}} := f\circ \pi_X \;=\; \rho\circ \pi_S,
    \]
    interpreted as equality in $\mathcal{C}_{\mathrm{stoch}}$.
    (If one adopts weaker, prior-dependent correctness, then $\pi_{\mathrm{out}}$ is best treated as a distinguished comparator kernel rather than literally equal to both composites.)

    \item[\textbf{Correctness relation.}]
    Strong (kernel) correctness is the commutativity relation
    \[
    \rho\circ \pi_S \;=\; f\circ \pi_X \qquad \text{as morphisms } E\leadsto Y.
    \]
    A weaker (prior-based) correctness statement is: for a chosen prior $\mu$ on $E$, the induced output laws agree:
    \[
    (\rho\circ \pi_S)_*\mu \;=\; (f\circ \pi_X)_*\mu \quad \text{in }\Prob(Y).
    \]

    \item[\textbf{$\pi_{X_H}:E\to X_H$.}]
    The honest-input projection from executions. In the explicit-randomness model take
    $\pi_{X_H}:=\mathrm{proj}_{X_H}\circ \pi_X$ (deterministic); otherwise define it as the corresponding projection on $E=\prod_i X_i$.

    \item[\textbf{$\mathrm{proj}_{X_H}:\prod_i X_i\to X_H$.}]
    Deterministic coordinate projection forgetting the adversarial coordinates (a deterministic kernel).

    \item[\textbf{$\pi_{\mathrm{full}}:E\leadsto B_{\mathrm{full}}$.}]
    Maximal transcript kernel capturing all data to be treated as observable in the ``full transcript'' model.

    \item[\textbf{$\mathrm{proj}_S:B_{\mathrm{full}}\to \prod_i S_i$.}]
    Deterministic measurable post-processing map extracting the share-vector component from the maximal transcript.

    \item[\textbf{$\mathrm{proj}_A:B_{\mathrm{full}}\to B_A$.}]
    Deterministic measurable post-processing map extracting the adversarially visible part of the maximal transcript.

    \item[\textbf{$\pi_A:E\leadsto B_A$.}]
    The adversary view kernel, required by the ambient-diagram relation to satisfy
    \[
    \pi_A \;=\; \mathrm{proj}_A\circ \pi_{\mathrm{full}}.
    \]

    \item[\textbf{$\iota_A:B_A\hookrightarrow B_{\mathrm{full}}$ (optional padding).}]
    If one wants a literal padding inclusion, one must choose a measurable ``default filler'' for the coordinates missing from $B_A$.
    Such a choice is extra data (not canonical in general); if provided, $\iota_A$ is interpreted as a deterministic measurable injection.
\end{description}

\section{A First Quantum Embedding Attempt (with a Quantum-Augmented Diagram)}
\label{sec:quantum}

This section sketches an interpretation of the ambient diagram in a quantum setting.
The key structural change is that the natural composition rule for ``systems side by side'' is tensor product, so the deterministic-product template must be re-read in a symmetric monoidal setting.

\subsection{Target category (Schr\"odinger picture for explicit arrows)}
We work in the category $\mathcal{C}_{\mathrm{q}}:=\mathbf{CPTP}_{\mathrm{fd}}$ whose objects are finite-dimensional Hilbert spaces and whose morphisms are completely positive trace-preserving (CPTP) maps between density operators.\footnote{One can equivalently work in the Heisenberg picture with completely positive unital maps between finite-dimensional C$^*$-algebras.}
The monoidal product is the tensor product $\otimes$.
Classical registers are modeled as Hilbert spaces with a preferred orthonormal basis; deterministic classical functions induce CPTP maps that push forward classical (diagonal) states.

\subsection{Why the diagram needs a quantum augmentation}
In genuinely quantum-adversary settings, the adversary's raw side information is typically a quantum system $Q_A$, and a classical transcript $B_A$ is obtained only after choosing a measurement/classicalization channel $M_A:Q_A\to B_A$.
Thus, if we want the diagram to reflect both possibilities (``keep the view quantum'' vs ``classicalize''), we extend the ambient diagram by adding $Q_A$ and $M_A$.

\begin{figure}[ht]
\centering
\begin{tikzcd}[row sep=large, column sep=large]
& E \arrow[r, "\pi_S"] \arrow[d, "\pi_X"'] \arrow[rr, bend left=20, "\pi_{\mathrm{full}}"] \arrow[dr, "\pi_{\mathrm{out}}"'] & \prod_i S_i \arrow[d, "\rho"] & B_{\mathrm{full}} \arrow[l, "\mathrm{proj}_S"'] \arrow[r, "\mathrm{proj}_{Q_A}"] & Q_A \arrow[d, "M_A"] \\
X_H & \prod_i X_i \arrow[l, "\mathrm{proj}_{X_H}"'] \arrow[r, "f"'] & Y & & B_A
\end{tikzcd}
\caption{Quantum-augmented ambient diagram: the adversary's raw view is a quantum system $Q_A$ obtained by a restriction $\mathrm{proj}_{Q_A}$ of the maximal transcript, and a chosen classicalization/measurement $M_A$ yields a classical transcript $B_A$ if desired.}
\label{fig:quantum-augmented-ambient}
\end{figure}

\subsection{Node-by-node interpretation (faithful list)}
\begin{description}
    \item[\textbf{$X_i$.}] Classical input registers, modeled as Hilbert spaces $H_{X_i}\cong \mathbb{C}^{|X_i|}$ with a preferred basis $\{\lvert x_i\rangle\}$.
    \item[\textbf{$\prod_i X_i$.}] The joint classical input register $H_{X}:=\bigotimes_i H_{X_i}$.
    \item[\textbf{$X_H$.}] Honest-input classical register $H_{X_H}:=\bigotimes_{i\in H} H_{X_i}$.
    \item[\textbf{$Y$.}] Classical output register $H_Y\cong \mathbb{C}^{|Y|}$ (preferred basis $\{\lvert y\rangle\}$).
    \item[\textbf{$S_i$ and $\prod_i S_i$.}] Transcript/share registers; in classical protocols these are classical registers. Write $H_S:=\bigotimes_i H_{S_i}$.
    \item[\textbf{$E$.}] Execution system. A typical model is
    \[
    H_E := H_X \otimes H_\Omega \otimes H_W,
    \]
    where $H_\Omega$ is a classical randomness register and $H_W$ is an internal quantum workspace (including any quantum communication transcripts before measurement).
    \item[\textbf{$B_{\mathrm{full}}$.}] Maximal transcript system; can be classical or quantum depending on modeling. In the augmented diagram it is a system from which both shares and the adversary view are obtained by post-processing.
    \item[\textbf{$Q_A$.}] Quantum adversarial side-information system (raw quantum view).
    \item[\textbf{$B_A$.}] Classical adversarial transcript system obtained after choosing a classicalization/measurement $M_A$ (optional; one can also keep the view quantum and omit $B_A$).
\end{description}

\subsection{Arrow-by-arrow interpretation (every named arrow)}
\begin{description}
    \item[\textbf{$\pi_X:E\to \prod_i X_i$.}]
    The CPTP map that discards internal registers and keeps the input register:
    \[
    \pi_X(\rho_E) := \mathrm{Tr}_{\Omega W}(\rho_E) \in \mathsf{D}(H_X).
    \]

    \item[\textbf{$\mathrm{proj}_{X_H}:\prod_i X_i\to X_H$.}]
    The CPTP map that discards adversarial input coordinates, i.e.\ partial trace over $X_A$:
    $\mathrm{proj}_{X_H}(\rho_X):=\mathrm{Tr}_{X_A}(\rho_X)$.

    \item[\textbf{$\pi_{X_H}:E\to X_H$.}]
    Defined by $\pi_{X_H}:=\mathrm{proj}_{X_H}\circ \pi_X$.

    \item[\textbf{$f:\prod_i X_i\to Y$.}]
    The ideal functionality as a \emph{classical deterministic channel}:
    on computational basis states, $\lvert x\rangle\mapsto \lvert f(x)\rangle$.
    Concretely, on diagonal density matrices (classical distributions) it pushes forward the law of $X$ to the law of $Y$.

    \item[\textbf{$\pi_S:E\to \prod_i S_i$.}]
    The protocol's transcript extraction channel: it outputs the (classical) transcript registers corresponding to the parties' local views.
    In a classical protocol this is induced by a deterministic function of $(x,\omega)$; in general it can be a CPTP channel from $H_E$ to $H_S$.

    \item[\textbf{$\rho:\prod_i S_i\to Y$.}]
    The reconstruction post-processing as a classical deterministic channel on transcript registers (CPTP map on commutative/diagonal states).

    \item[\textbf{$\pi_{\mathrm{out}}:E\to Y$.}]
    The output-observation channel. In an idealized perfect-correctness model one enforces the equality of CPTP maps
    \[
    \pi_{\mathrm{out}} = f\circ \pi_X = \rho\circ \pi_S.
    \]

    \item[\textbf{$\pi_{\mathrm{full}}:E\to B_{\mathrm{full}}$.}]
    A maximal transcript channel, bundling all records of interest (classical and/or quantum) into a single system.

    \item[\textbf{$\mathrm{proj}_S:B_{\mathrm{full}}\to \prod_i S_i$.}]
    Post-processing that extracts the share/transcript subsystem from the maximal transcript (a discard of the complementary subsystem).

    \item[\textbf{$\mathrm{proj}_{Q_A}:B_{\mathrm{full}}\to Q_A$.}]
    Post-processing that extracts the adversary's quantum side-information subsystem (again, a discard/partial trace).

    \item[\textbf{$M_A:Q_A\to B_A$ (classicalization/measurement).}]
    A chosen CPTP measurement/instrument that outputs a classical register $B_A$.
    This is extra modeling data; different choices of $M_A$ correspond to different ``what the adversary reads out'' assumptions.

    \item[\textbf{$\pi_A:E\to B_A$.}]
    In the classicalized model, define
    \[
    \pi_A := M_A\circ \mathrm{proj}_{Q_A}\circ \pi_{\mathrm{full}}.
    \]
    If one keeps the adversary view quantum, the analogue is the channel $E\to Q_A$ given by $\mathrm{proj}_{Q_A}\circ \pi_{\mathrm{full}}$.

    \item[\textbf{$\mathrm{proj}_A:B_{\mathrm{full}}\to B_A$ (classicalized).}]
    If one wants to keep the original ambient-diagram arrow name $\mathrm{proj}_A$, set
    \[
    \mathrm{proj}_A := M_A\circ \mathrm{proj}_{Q_A}.
    \]

    \item[\textbf{$\iota_A:B_A\hookrightarrow B_{\mathrm{full}}$ (padding, optional).}]
    A ``padding inclusion'' requires choosing an ancilla state $\sigma$ for the complementary subsystem and defining a state-preparation channel
    \[
    \iota_A(\rho_{B_A}) := \rho_{B_A}\otimes \sigma.
    \]
    This is non-canonical (depends on $\sigma$) and should be treated as bookkeeping rather than an operational adversary action.
\end{description}

\begin{remark}[Not claimed]
This section does not claim a composable quantum security theorem; it only records how the ambient-diagram \emph{syntax} can be interpreted in a standard quantum channel category, and where the cartesian-product intuition must be replaced by monoidal/discard structure.
\end{remark}

\section{Embedding into \texorpdfstring{$\mathbf{Top}$}{Top}}
\label{sec:top}

To interpret the ambient diagram in $\mathbf{Top}$, one must choose topological spaces for all nodes and require each arrow to be continuous.

\subsection{Node-by-node interpretation (faithful list)}
\begin{description}
    \item[\textbf{$X_i$.}] Topological input spaces (chosen by the model).
    \item[\textbf{$\prod_i X_i$.}] Product space with the product topology.
    \item[\textbf{$X_H$.}] Honest-input product subspace $\prod_{i\in H}X_i$.
    \item[\textbf{$S_i$ and $\prod_i S_i$.}] Topological share/transcript spaces and their product.
    \item[\textbf{$Y$.}] Topological output space.
    \item[\textbf{$\Omega$.}] Topological internal randomness/state space (if modeled explicitly).
    \item[\textbf{$E$.}] Topological execution space (often $E\cong (\prod_i X_i)\times \Omega$).
    \item[\textbf{$B_{\mathrm{full}}$.}] Topological maximal-transcript space.
    \item[\textbf{$B_A$.}] Topological adversary-view space.
\end{description}

\subsection{Arrow-by-arrow interpretation (every arrow in Figure~\ref{fig:master-diagram})}
\begin{description}
    \item[\textbf{$\pi_X:E\to \prod_i X_i$.}] Continuous projection/observation forgetting internal state (e.g.\ projection if $E=(\prod_i X_i)\times\Omega$).
    \item[\textbf{$f:\prod_i X_i\to Y$.}] Continuous map (ideal functionality) in the topological model.
    \item[\textbf{$\pi_S:E\to \prod_i S_i$.}] Continuous protocol transcript map.
    \item[\textbf{$\rho:\prod_i S_i\to Y$.}] Continuous reconstruction map.
    \item[\textbf{$\pi_{\mathrm{out}}:E\to Y$.}] Continuous output observation; the diagram requires $\pi_{\mathrm{out}}=f\circ \pi_X=\rho\circ \pi_S$.
    \item[\textbf{$\mathrm{proj}_{X_H}:\prod_i X_i\to X_H$.}] Continuous coordinate projection.
    \item[\textbf{$\pi_{X_H}:E\to X_H$.}] Continuous map, typically $\pi_{X_H}=\mathrm{proj}_{X_H}\circ \pi_X$.
    \item[\textbf{$\pi_{\mathrm{full}}:E\to B_{\mathrm{full}}$.}] Continuous maximal transcript map.
    \item[\textbf{$\mathrm{proj}_S:B_{\mathrm{full}}\to \prod_i S_i$.}] Continuous post-processing/forgetful map.
    \item[\textbf{$\mathrm{proj}_A:B_{\mathrm{full}}\to B_A$.}] Continuous post-processing/forgetful map to adversarial view.
    \item[\textbf{$\pi_A:E\to B_A$.}] Continuous adversary view map satisfying $\pi_A=\mathrm{proj}_A\circ \pi_{\mathrm{full}}$.
    \item[\textbf{$\iota_A:B_A\hookrightarrow B_{\mathrm{full}}$ (optional padding).}] A continuous inclusion exists only after choosing a continuous ``padding''/section datum; otherwise omit it as non-canonical bookkeeping.
\end{description}

\begin{remark}[Research direction]
Nontrivial topological content typically requires choosing topologies that reflect an actual geometry (parameter spaces, asymptotic limits, quotient topologies from encodings, etc.). The ambient-diagram template then becomes a bridge between protocol structure and homotopy-theoretic notions (e.g.\ whether certain post-processing maps behave like fibrations/cofibrations under a chosen model structure).
\end{remark}

\section{Embedding into \texorpdfstring{$\mathbf{Man}$}{Man}}
\label{sec:man}

To interpret the ambient diagram in $\mathbf{Man}$ (smooth manifolds), one must model each node as a manifold and each arrow as a smooth map.

\subsection{Node-by-node interpretation (faithful list)}
\begin{description}
    \item[\textbf{$X_i$.}] Smooth input manifolds (often approximation/relaxation of discrete inputs).
    \item[\textbf{$\prod_i X_i$.}] Product manifold (when it exists in the chosen manifold category; typically true for finite products).
    \item[\textbf{$X_H$.}] Honest-input product manifold $\prod_{i\in H}X_i$.
    \item[\textbf{$S_i$ and $\prod_i S_i$.}] Smooth transcript manifolds and their product.
    \item[\textbf{$Y$.}] Smooth output manifold.
    \item[\textbf{$\Omega$.}] Smooth internal randomness/state manifold (in an approximation regime).
    \item[\textbf{$E$.}] Smooth execution manifold (often $E\cong (\prod_i X_i)\times \Omega$).
    \item[\textbf{$B_{\mathrm{full}}$.}] Smooth maximal-transcript manifold (requires a modeling choice/relaxation).
    \item[\textbf{$B_A$.}] Smooth adversary-view manifold (requires a modeling choice/relaxation).
\end{description}

\subsection{Arrow-by-arrow interpretation (every arrow in Figure~\ref{fig:master-diagram})}
\begin{description}
    \item[\textbf{$\pi_X:E\to \prod_i X_i$.}] Smooth projection/observation forgetting internal state (e.g.\ projection if $E=(\prod_i X_i)\times\Omega$).
    \item[\textbf{$f:\prod_i X_i\to Y$.}] Smooth functionality map (in the relaxed model).
    \item[\textbf{$\pi_S:E\to \prod_i S_i$.}] Smooth protocol transcript map (in the relaxed model).
    \item[\textbf{$\rho:\prod_i S_i\to Y$.}] Smooth reconstruction map.
    \item[\textbf{$\pi_{\mathrm{out}}:E\to Y$.}] Smooth output observation with $\pi_{\mathrm{out}}=f\circ\pi_X=\rho\circ\pi_S$ enforced as an identity of smooth maps.
    \item[\textbf{$\mathrm{proj}_{X_H}:\prod_i X_i\to X_H$.}] Smooth coordinate projection.
    \item[\textbf{$\pi_{X_H}:E\to X_H$.}] Smooth map, typically $\pi_{X_H}=\mathrm{proj}_{X_H}\circ \pi_X$.
    \item[\textbf{$\pi_{\mathrm{full}}:E\to B_{\mathrm{full}}$.}] Smooth maximal transcript map (depends on the transcript relaxation).
    \item[\textbf{$\mathrm{proj}_S:B_{\mathrm{full}}\to \prod_i S_i$.}] Smooth post-processing/forgetful map.
    \item[\textbf{$\mathrm{proj}_A:B_{\mathrm{full}}\to B_A$.}] Smooth post-processing map to the adversary-view manifold.
    \item[\textbf{$\pi_A:E\to B_A$.}] Smooth adversary view map satisfying $\pi_A=\mathrm{proj}_A\circ \pi_{\mathrm{full}}$.
    \item[\textbf{$\iota_A:B_A\hookrightarrow B_{\mathrm{full}}$ (optional padding).}] Requires a smooth padding choice/section; otherwise omit as non-canonical bookkeeping.
\end{description}

\begin{remark}[Not claimed]
Embedding into $\mathbf{Man}$ is typically an \emph{approximation} program: discrete transcript alphabets are replaced by smooth parameterizations (e.g.\ simplices, Euclidean embeddings, statistical manifolds). The purpose is to enable differential-geometric tools (local sensitivity, Fisher metrics, smooth projection principles), not to assert that the original protocol literally lives in $\mathbf{Man}$.
\end{remark}

\end{document}

